\section{Lifting a coordinate on a mixed ramified letter}

% BEGIN X81M1
\paragraph{X81M1. The two branch values and the cost of a down-branch constant.}
On a mixed letter over $d$ denote the up deep endpoint by $b$,
the down endpoint by $b_D$, and put $B=b_D+S$ for the integral
width $S$ of H80M2. Thus $\eta<b,b_D$ and $b,B\le h$.
The dual sign reversal gives endpoints $s-1$ and $s+1$ respectively
in H80M2. The two values of $s$ use the archimedean data of their
respective embeddings. We allow both possible orders of $b,B$.

Write $a_U,a_D\in O$ for the up and down branch values of $v$.
At a fixed $K_0$-embedding the roots giving one label form one
orbit of $\mathrm{Gal}(K_1/K)$, by the extension-of-embeddings
argument in W79H2. The two values are distinct and
$v_p(a_D-a_U)=1/2$: by A81R3 they each have valuation $1/2$ and
their ratio reduces to the nontrivial square-class sign $-1$,
distinct from $1$ for odd $p$. Put $D_U=v-a_U$, $D_D=v-a_D$,
$E=D_U D_D$. At each branch $D_U$ or $D_D$ respectively is a
formal local coordinate over $L$ by the simple zero in P81F2.
An integral series in $v$ has integral coefficients in its
Taylor expansion in either of these coordinates since the two
$a$'s belong to the maximal ideal.

On this letter an $\mathrm{IsCoord}$ is $f=A/E^h$, $A\in O[[v]]$.
At the up branch $A$ vanishes to order at least $h-b$; at the
down branch $A-x_0E^h$ vanishes to order at least $h+B$ for
some $x_0\in L$. Divide $A$ integrally by $D_U^{h-b}$
by successive division at $a_U$, obtaining $A'$. At the down
branch $D_U$ is a rational Taylor unit. With
$\delta=a_D-a_U$ the precise congruence is
\[
      A'\equiv x_0 D_D^h(D_D+\delta)^b\pmod{D_D^{h+B}}.
\]
If $b<B$, its coefficient at degree $h+b$ gives $x_0\in O$.
If $B\le b\le B+2\eta$, its coefficient at degree $h+B-1$
gives $x_0\binom b{B-1}\delta^{b-B+1}\in O$.
By Legendre's factorial formula the valuation of a nonzero
binomial $\binom b k$ is at most $\lfloor\log_p b\rfloor$:
each nonzero level in the difference of the three floor sums
contributes at most one. Hence in this range
$p^{y_0}x_0\in O$, where
$y_0=\lceil\eta+1/2+\log_p h\rceil$, as in A81R4.
Both congruences hold at all primitive roots inducing the same
label, since the Taylor-coordinate changes are invertible and
the single rational quotient $A/E^h$ is the same.
% END X81M1

% BEGIN X81M2
\paragraph{X81M2. Lifting in the close-branch range.}
Use the frame of P81F1 and following claims with its model
compatibility V81P2. Denote the integral constant module on
the predecessor by $D_j=N_{c-1}(T_j)^\Delta/vN_{c-1}(T_j)^\Delta$.
The period map of P81F1 embeds its rationalization into the
crystalline space on the current letter by Frobenius of the
preceding constant term. It is the same map for each primitive
root, by W79H3. Write $F_D:D_1\to D_2$ for the induced map.
For the tested predecessor vector $x$ write $d\in D_2$ for its
class. Its up value lies in $\mathrm{Fil}^{b}$. In either case
of X81M1 choose $y\le y_0$ so $o=p^y x_0\in O$.

The model tensor $(p^yf-o)d$ belongs to the target $\mathcal N_2$.
Indeed the numerator $p^y A-oE^h$ vanishes to order at least
$h-b$ at the up branch and to order at least $h+B$ at the down
branch. On the up branch every possible degree $i$ therefore
satisfies $h-i\le b$, giving the required filtration test. On the
down branch $h-i\le-B$ and the full period space lies there
since $-B$ is a lower bound on jumps. These arguments use P81F1
and the integrality of $o$.

We record separately that $fd$ is invariant modulo $v\mathcal N_2$,
without spending $p^y$ in this assertion. In fact
$g=E^h(\gamma f-f)/v$ is integral by the last $\mathrm{IsCoord}$
clause. The substitution fixes each branch. At the up branch
$f$ has order at least $-b$ and at the down branch $f-x_0$ has
order at least $B$, with $\gamma x_0=x_0$. Since $v$ is nonzero
at both branches, the orders of $g$ are at least $h-b$ and $h+B$
there, respectively. The same two filtration tests give
$gd/E^h\in\mathcal N_2$. This argument also applies to any
source vector whose up value belongs to $\mathrm{Fil}^b$.

Let $w$ represent $\alpha(fd)$ as in the tested statement; it is
invariant modulo $vN_2$ by V81P2. Choose
$n\in N_c(T_2)^\Delta$ representing $\alpha((p^yf-o)d)$ using
V81P2. Lift $p^\rho n$ to $n'\in N_c(T_1)^\Delta$ by A81R9.
These vectors satisfy the $\mathrm{InvS}$ conditions modulo $v$
by $\mathrm{NinLam}$ and $\mathrm{GamPi}$ in A81R7.
For the constant tensor $d$ use A81R9 on the preceding letter
to choose $d'\in D_1$ above $p^\rho d$. A constant model tensor
belongs to $\mathrm{InvL}$ because its difference is zero.
Choose representatives $z,z'$ of $\alpha(d),\alpha(d')$.
Both lie in the required $\mathrm{InvS}$ by V81P2.
Naturality gives $Fz'-p^\rho z\in\mathrm{JP}_2$, and linearity
gives $p^y w-n-oz\in\mathrm{JP}_2$. Here
$\mathrm{JP}_2\subseteq v^M N(T_2)^\Delta$: multiply its
Frobenius-span factor by $E^h$ to make it integral by W79H4,
and its defining multiplier is $v^ME^{3h}$. Thus
$n'+oz'$ lifts $p^{\rho+y}w$ modulo $vN(T_2)^\Delta$, as required.
% END X81M2

% BEGIN X81M3
\paragraph{X81M3. The separated branches give a filtered source lift.}
Suppose now $b>B+2\eta$. Put
$q=D_U^{h-b}D_D^{h+B}$ and $f_0=q/E^h$.
For the tested $d\in D_2$ of X81M2,
$f_0d\in\mathcal N_2$ by the same Taylor tests: the numerator
has up order $h-b$ and down order $h+B$. Use V81P2 to choose
$n\in N_c(T_2)^\Delta$ representing $\alpha(f_0d)$.
Choose $n'\in N_c(T_1)^\Delta$ above $p^\rho n$ by A81R9,
and by V81P2 choose $\nu\in\mathcal N_1$ representing
$\beta(n')$. The maps of the frame are functorial for $F$ and
$\beta\alpha=p^{2\sigma}$ in their quotients. Therefore
$F_D\nu-p^{\rho+2\sigma}f_0d$ belongs to $\mathrm{JLA}_2$,
with $F_D$ extended to coefficients. Multiplying by $E^h$ gives
an integral series $H=E^h\nu\in D_1[[v]]$ satisfying
\[
 F_D H\equiv C q d\pmod {v^M E^{4h} D_2[[v]]},
                 \qquad C=p^{\rho+2\sigma}.
\]
This uses the exact quotient defining the frame, with its factor
$E^{3h}$ before clearing the denominator.

Because $\nu\in\mathcal N_1$, its numerator coefficient of
degree $i$ lies in $\mathrm{Fil}_D^{h-i}$ at the down branch
and in $\mathrm{Fil}_U^{h-i}$ at the up branch. The assertion
holds in the coordinates $D_D,D_U$ of X81M1: an invertible
change of local variable with zero constant mixes a coefficient
only with earlier degrees, whose filtration conditions are
stronger. On $D_1\otimes L$ the down filtration above $\eta$
vanishes, since all source jumps there are central or negative
deep. Thus $H$ vanishes to order $h-\eta$ at the down branch.
By integral division in a basis of $D_1$ write
$H=D_D^{h-\eta}H_1$ with $H_1\in D_1[[v]]$.
At the up branch $D_D$ is a rational Taylor unit; division by
its power preserves the same coefficient filtration bounds.
The congruence there now gives, to degree strictly below $4h$,
\[
    F_D H_1\equiv C D_U^{h-b}(D_U+a_U-a_D)^{B+\eta}d
                    \pmod {D_U^{4h}}.
\]
Take the coefficient of $D_U^j$ for
$j=h-b+B+\eta<h-\eta$. It is an integral vector $d'\in D_1$,
and its image is $Cd$ by the monic highest coefficient on the
right. The up filtration of $d'$ is at least $h-j>\eta$.
The source has no jumps strictly between $\eta$ and $b$, so
this is $d'\in\mathrm{Fil}_U^b$. This argument uses the common
rational crystalline space for both branch tests, with the two
separate labelled filtrations, as in X81M2. In particular the
down vanishing above is a statement in that space and follows
from injectivity of its period-component map.

Choose a predecessor Wach representative for $d'$. The final
$\mathrm{IsCoord}$ argument in X81M2 applies to $fd'$:
the up vector now has filtration at least $b$, and the down
filtration again contains the full source space at $-B$.
Thus $fd'\in\mathrm{InvL}$ on the source and its forward frame
representative lies in $\mathrm{InvS}$ modulo $v$ by V81P2.
Frame naturality and $F_Dd'=Cd$ show that this representative
maps to $Cw$ modulo $\mathrm{JP}_2\subset vN(T_2)^\Delta$.
Combining with X81M2 and scaling either lift by the unused
nonnegative power yields the original mixed
$\mathrm{ThirdPropositionSlot}$ at
$\rho+\max(2\sigma,y_0)$ for all the row maps at $d$.
% END X81M3
